\documentclass[11pt,largemargins, letter]{homework}

\newcommand{\hwname}{Tobias Jakobi}
\newcommand{\hwemail}{tjakobi@arizona.edu}
\newcommand{\hwtype}{Project}
\newcommand{\hwnum}{1}
\newcommand{\hwclass}{CTS505}
\newcommand{\hwlecture}{2026}
\newcommand{\hwsection}{-11}

\usepackage[margin=1in]{geometry}
\usepackage{fontawesome5}

\renewcommand{\questiontype}{Task}

\begin{document}
\maketitle

The solutions for the tasks listed below should be saved in separate \texttt{.R} files, e.g. \texttt{task1.R}, \texttt{task2.R}.

Please include loaded libraries, loading of files, and all other commands in the files.

\question

 \begin{alphaparts}
   \questionpart Download a list of gene IDs with expression values and false discovery rate (FDR) columns from \texttt{https://links.jakobilab.org/cts505\_genes}.
   \questionpart Identify the organism used for the experiment based on the provided gene IDs for the next tasks.
 \end{alphaparts}

\question

 \begin{alphaparts}
   \questionpart Read the provided input file into R with an appropriate function.
   \questionpart Install and use the appropriate \texttt{org.*.db} annotation package from Bioconductor that matches the identified organism (e.g., \texttt{org.Hs.eg.db} for human, \texttt{org.Mm.eg.db} for mouse) to retrieve gene names and Uniprot IDs for the provided gene IDs.
   \questionpart Print the top 10 highest expressed genes based on the logFC value into a new CSV file, \texttt{top10.csv}.
   Include the gene names and Uniprot IDs.
 \end{alphaparts}

\question

 \begin{alphaparts}
   \questionpart Use the \texttt{ggplot2} package to create an X-Y graph that plots \texttt{log10(FDR value)} on the Y axis and the logFC value on the X axis.
   \questionpart Set a title and descriptive X and Y axis labels and save the plot to PDF file named \texttt{xy\_plot.pdf}.
 \end{alphaparts}

\pagebreak
 \question

\begin{itemize}

   \item Download the list of gene counts stored in the following file \\\texttt{https://links.jakobilab.org/p2\_counts}
   \item Identify the organism used for the experiment based on the provided gene IDs for the next tasks.
\end{itemize}

 \begin{alphaparts}

   \questionpart Establish a reasonable grouping of the provided samples based on the sample names.
   \questionpart Read the provided input file into \faRProject \hspace{0.1mm} with an appropriate function.
   \questionpart Install and use the appropriate \texttt{org.*.db} annotation package from Bioconductor for the identified organism to retrieve gene names for the gene IDs.
 \end{alphaparts}

\question

 \begin{alphaparts}
   \questionpart Develop a suitable design formula as discussed in the class to perform differential gene expression analysis between your chosen groups.
   \questionpart Use the \texttt{edgeR} workflow to build an \texttt{DGElist} object.
   \questionpart Using \texttt{edgeR}, perform filtering and normalization of the imported count data.
   \questionpart Lastly, perform differential gene expression analysis and note the number of up- and down regulated genes.
 \end{alphaparts}

\question

 \begin{alphaparts}
   \questionpart Save a result table with differentially expressed genes filtered by p-value and FDR (0.05) and sorted by logFC.
 \end{alphaparts}


Please send your R scripts \& plots as attachment via email to \texttt{tjakobi@arizona.edu}.

In case of questions please do not hesitate to contact me.

\end{document}
